\chapter*{Individual Contributions}
\addcontentsline{toc}{backmatter}{Individual Contributions}
\label{cp:Contributions}
\markboth{\textit{Individual Contributions}}{\textit{Individual Contributions}}

\parindent0pt

The project was designed together in weekly team meetings. Each member then owned the parts listed below, and all
members reviewed each other's code and contributed to testing, the presentation and this report.

\vspace{0.4em}
\renewcommand{\arraystretch}{1.25}
\begin{tabularx}{\linewidth}{@{}>{\raggedright\arraybackslash}p{3.6cm} X@{}}
\toprule
\textbf{Member} & \textbf{Individual contribution} \\
\midrule
\textbf{Rithvik Jami}\newline {\small\itshape Team Leader}\newline {\scriptsize CS23BTECH11022} &
\textbf{Simulation core, radio network and integration.} Designed the simulator architecture and the per-tick timing
model (beacons, sensing and ranging at the same tick); implemented the configuration system with validation, the
paired random streams and the pairing check; built the radio layer (disk graph with packet loss, gossip relaying with
priorities, re-query); wrote the experiment runner and sweep infrastructure; ran all experiments on the GPU server; led
the first-build theory studies ($\lambda_2$, Erd\H{o}s--R\'enyi robustness, degradation); coordinated the team and the
integration of all modules. \\
\midrule
\textbf{Praneeth Chamarthy}\newline {\scriptsize CS23BTECH11012} &
\textbf{Mission agent, mapping and planning.} Implemented the drone's mission state machine (transit, explore, confirm,
audit, return) and the W-MSR formation in transit; the interval terrain map and the provenance-tagged coverage maps;
the navigation grid, Dijkstra path planning and the decentralised geodesic Voronoi partition with work stealing;
target selection, hysteresis and look-before-move; the dead-lock fixes (camera-located liar, escape and
route-around) and the fair \tcode{trust\_gate} comparison; compiled and edited this report. \\
\midrule
\textbf{Muvvala Sai Ram Suhas}\newline {\scriptsize CS23BTECH11038} &
\textbf{Epistemic layer.} Implemented the S5 possible-worlds hypotheses over liar sets; the proof rules R1--R6 and
their soundness arguments; the belief and knowledge standards over provenance bitmasks; threshold handling of
accusations; uniformly sampled audits; the knowledge gate $\Kn_i\mathit{Safe}$ with hold, re-query and right of way;
nested knowledge and the cross-check against the Kripke-model engine; the soundness-under-attack tests. \\
\midrule
\textbf{Jakkala Naga Rohith Kumar}\newline {\scriptsize CS23BTECH11021} &
\textbf{World, sensing and visualisation.} Built the city generator with random collapse and the scenarios; the camera
model with ray casting, the single \tcode{viewed} function for coverage and detection, and drone sightings; the
optional CUDA ray-casting path verified against NumPy; the Three.js 3D replay viewer and its exporter; the MP4 video
export. \\
\midrule
\textbf{Gona Sanjana}\newline {\scriptsize CS23BTECH11019} &
\textbf{Adversary model and evaluation.} Implemented the liar agent (motions, range-aware and other position lies, fake
coverage); the W-MSR aggregators and the $(r,s)$-robustness checker; the SPRT alive/dead belief and the survivor
localisation study; the mission metrics and STL specifications; the statistical analysis (confidence intervals, paired
Wilcoxon tests) and the result figures. \\
\midrule
\textbf{Prabhu Teja}\newline {\scriptsize CS23BTECH11050} &
\textbf{Demo and documentation support.} Set up the demonstration scenario and seeds, tested the Streamlit dashboard,
prepared slide figures, and helped write the presentation and the plain-language project guide. \\
\bottomrule
\end{tabularx}
